% ==========================================
% RESUMEN
% ==========================================

\chapter*{Resumen}

Este reporte presenta el desarrollo de una práctica de cómputo paralelo en lenguaje C orientada a la creación y gestión de procesos hijo para resolver tareas de manera concurrente. En la introducción se contextualiza la importancia del paralelismo y se plantean los objetivos y alcances del trabajo. En el marco teórico se describen los tipos de datos en C, su representación en memoria y su impacto en la programación correcta. Posteriormente, se estudian los procesos en sistemas tipo Unix, incluyendo su ciclo de vida, llamadas al sistema como \texttt{fork()}, \texttt{exec()} y \texttt{wait()}, así como casos de procesos padre, hijo, zombi y huérfano. También se aborda el manejo de archivos en C y su relevancia para lectura, escritura y persistencia de resultados en contextos concurrentes. En el capítulo de programa se documenta la solución implementada para procesar un archivo CSV, asignar trabajo por columna a procesos hijo y consolidar salidas en un archivo común con validaciones y sincronización. Finalmente, las conclusiones integran aprendizajes técnicos, dificultades encontradas, aplicaciones prácticas y propuestas de mejora.